\documentclass[10pt]{article}

\usepackage[utf8]{inputenc}

\usepackage[T1]{fontenc}

\usepackage{amsmath}

\usepackage{amsfonts}

\usepackage{amssymb}

\usepackage[version=4]{mhchem}

\usepackage{stmaryrd}

\usepackage{mathrsfs}

\usepackage{graphicx}

\usepackage[export]{adjustbox}

\graphicspath{ {./images/} }

\usepackage{bbold}



\begin{document}

ХЕЛЛИ ТЕОРЕМА - 1) Х. т. о п е р е с е ч е н и и в Ы пусть $K$ - семейство из не менее чем $n+1$ выпуклых множеств в $n$-мерном афффинном пространстве $A^{n}$, причем $K$ - конечно или каждое множество из $K$ - компактно; тогда, если каждые $n+1$ из множеств семейства имеют общую точку, то существует точка, общая всем множествам семейства $K$.



X. т. посвящены многие исследования, относящиеся к ее приложениям, доказательству различных аналогов и предложений типа X. т., ее обобщений, напр. в вопросах чебышевского приближения, в решениях освещения задач, в теории выпуклых тел. Часто Х. т. фигурирует в доказательствах комбинаторных утверждений следующего типа: если в нек-ром семействе каждое подсемейство из $k$ членов обладает определенным свойством, то этим свойством обладает и все семейство. Напр., если $a$ и $b$ - две точки множества $K \subset A^{n}$, то выражение «a видно из $b$ в $K$ » обозначает, что отрезок $[a, b]$ принадлежит $K$. Пусть компакт $K \subset A^{n}$ обладает свойством, тто для каждых $n+1$ точек из $K$ существует точка в $K$, из к-роӥ видны эти точки, тогда в $K$ существует точка, из к-рой видны все точки $K$, т. е. $K-$ 3 в е 3 д н о е м н $о$ же с т в 0 .



Большинство аналогов $X$. т. и ее обобщений связаны с различными вариантами гонятия «выпуклосты». Напр., пусть $S^{n}$ - евклидова сфера; множество наз. в ы с каждой пароӥ диаметрально непротивоположных точек ого содержит соединяющую әти точки менышую дугу определяемой ими большой окружности. Если семсйство выпуклых по Робинсону замкнутых множеств шространства $S^{n}$ таково, что каждые $2(n+1)$ его элементов обладают непустым пересечением, то и все элементы этого семеӥства обладают непустым пересечением.



Х. т. установлена Э. Хелли (E. Helly, 1913).



Jum.: [1] Д а н ц е р Ј., $\quad \Gamma \mathrm{p}$ ю 6 а ум Б., К К и В., Теорема Хелли, пер. с англі, М., 1968. П. С. Солтан.



\begin{enumerate}

  \setcounter{enumi}{1}

  \item Х. т. в т е о р и и функ ци й: если последовательность функций $g_{n}, n=1,2, \ldots$, с ограниченной на отрезке $[a, b]$ вариацией сходится в каждой точке этого отрезка к пек-рой функции $g$, причем вариации $\underset{a}{\vee} g_{n}$ всех функций $g_{n}$ ограничены в совокушности:

\end{enumerate}



\[

\stackrel{b}{\vee} g_{n} \leqslant c, \quad n=1,2, \ldots,

\]



то предельная функция $g$ также имеет ограниченную на отрезке $[a, b]$ вариацию и для любой непрерывной на $[a, b]$ функции $f$ справедливо равенство



\[

\lim _{n \rightarrow \infty} \int_{a}^{b} f(x) d g_{n}(x)=\int_{a}^{b} f(x) d g(x) .

\]



Л. Д.. Кудрявчев.



ХЕЛЛИНГЕРА ИНТЕГРАЛ - интеграл тіпа Римана от функции множества $f(E)$. Если $(X, \mu)$-пространство с конечной неотрицательной несингулярной мерой, $f(E), \quad E \subset X,-$ віолне аддитивная функция, причем $f(E)=0$, если $\mu E=0 ; \delta=\left\{E_{n}\right\}_{1}^{N}$ - разбиение $X$, то



\[

S_{\delta}=\sum_{n=1}^{N} \frac{f^{2}\left(E_{n}\right)}{\mu E_{n}}

\]



и интегралом Хеллингера функции $f(E)$ по $X$ наз.



\[

\int_{X} \frac{f^{2}(d E)}{d \mu}=\sup _{\delta} S_{\delta}

\]



если он конечен. X. и. можно рассматривать и как предел по направленному множеству разбиений: $\delta_{1}<\delta_{2}$, если $\delta_{2}$ есть подразбиение $\delta_{1}$. Если существует суммируемая фуннция $\varphi(x): X \rightarrow R$ такая, что $f(E)$ есть интеграл Лебега $\int_{E} \varphi d \mu$, то Х. и. выражается через интеграл Лебега



\[

\int_{X} \frac{f^{2}(d E)}{d \mu}=\int_{X} \varphi^{2} d \mu .

\]



Э. Хеллингер [1] дал определение интеграла для $X=[a, b]$ в терминах функций точки.



Jum.: [1] H e l l i i g g e r E., «J. reine und angew. Math.», 1909\$, Вd 136, S. $210-71 ; .2]$ С м и р і о в В. И., Курс высшей

математики, т. 5 , М., 1959.

И. А. Виноградова.



ХЕЛЛИНГЕРА РАССТОЯНИЕ - расстояние между вероятностными мерами, выраженное в терминах Xeллингера интеграла. Пусть на измеримом пространстве $(\mathfrak{X}, \mathscr{B})$ задано семейство вероятностных мер $\left\{P_{\theta}\right\}, \theta \in \Theta$, абсолютно непрерывных относительно нек-рой б-конечной меры $\mu$ на $\mathscr{B}$.



X. p. между мерами $P_{\theta_{1}}$ и $P_{\theta_{2}}\left(\theta_{1}, \theta_{2} \in \Theta\right)$ олределяется по формуле



\[

\begin{aligned}

& r\left(\theta_{1}, \theta_{2}\right)=\left\{2\left[1-H\left(\theta_{1}, \theta_{2}\right)\right]\right\}^{1 / 2}= \\

= & \left\{\int_{\Re}\left[\sqrt{\frac{d^{\mathrm{P}} \theta_{1}}{d \mu}}-\sqrt{\frac{d^{\mathrm{P}} \theta_{2}}{d \mu}}\right]^{2} d \mu\right\}^{1 / 2},

\end{aligned}

\]



где



\[

H\left(\theta_{1}, \theta_{2}\right)=\int_{\mathfrak{X}} \sqrt{\frac{d \mathrm{P}_{\theta_{1}}}{d \mu}} \sqrt{\frac{d \mathrm{P}_{\theta_{2}}}{d \mu}} d \mu

\]



\begin{itemize}

  \item интеграл Хеллингера. X. р. не зависит от выбора меры $\mu$ и обладает следующими свойствами:

\end{itemize}



\begin{enumerate}

  \item $0 \leqslant r\left(\theta_{1}, \theta_{2}\right) \leqslant \sqrt{2}$;



  \item $r\left(\theta_{1}, \theta_{2}\right)=\sqrt{2}$ тогда и только тогда, гогда меры $\mathrm{P}_{\theta_{1}}$ и $\mathrm{P}_{\theta_{2}}$ сингулярны;



  \item $r\left(\theta_{1}, \theta_{2}\right)=0$ тогда и только тогда, когда $\mathrm{P}_{\theta_{1}}=\mathrm{P}_{\theta_{2}}$.



\end{enumerate}



Пусть



\[

\begin{gathered}

\left\|\mathrm{P}_{\theta_{1}}-\mathrm{P}_{\theta_{2}}\right\|=\sup _{B \in \mathscr{B}}\left|\mathrm{P}_{\theta_{1}}(B)-\mathrm{P}_{\theta_{2}}(B)\right|= \\

=\frac{1}{2} \int_{\mathfrak{X}}\left|\frac{d \mathrm{P}_{\theta_{1}}}{d \mu}-\frac{d \mathrm{P}_{\theta_{2}}}{d \mu}\right| d \mu

\end{gathered}

\]



\begin{itemize}

  \item расстояние по вариации между мерами $\mathrm{P}_{0_{1}}$ и $\mathrm{P}_{\theta_{2}}$. Тогда

\end{itemize}



\[

\frac{1}{2} r^{2}\left(\theta_{1}, \theta_{2}\right) \leqslant\left\|\mathbf{P}_{\theta_{1}}-\mathbf{P}_{\theta_{2}}\right\| \leqslant r\left(\theta_{1}, \theta_{2}\right) .

\]



Jіит.: [1] Г о Х.-С., Гауссовские меры в банаховых пространствах, пер. с англ., М., 1979; [2] К р а м е р Г., Математические методы статистики, пер. с англ., 2 изд., М., 1975; [3] И б $\mathrm{p}$ а г им о в И. А., Х а с ь м и н с К и й Р. З., Асимптотическая

теория оценивания, М., 1979; [4] 3 о ло т а е в В. М., «Зап.



\begin{center}

\includegraphics[max width=\textwidth]{2023_07_09_a01006407529fd695617g-1(1)}

\end{center}



ХЕФЛИГЕРА СТРУКТУРА коразмерности $q$ и класса $C^{r}$ на топологич. пространстве $X$-структура, огределяемая с помощью х е ф лг и г е р о в с к о г о а т л а с а (наз. также хеф ли г е р о в с к им ко ци к ло м) $\left\{U_{\alpha}, \Psi_{\alpha \beta}\right\}$, где $U_{\alpha}$-открытые подмножсства, покрывающие $X$, а



\[

\Psi_{\alpha \beta}: U_{\alpha} \cap U_{\beta} \rightarrow \Gamma_{r}^{q}, \quad u \longrightarrow \Psi_{\alpha \beta, u}

\]



\begin{itemize}

  \item нешрерывные отображения $U_{\alpha} \cap U_{\beta}$ в пучок $\Gamma_{r}^{q}$ ростков локальных $C^{r}$-диффеоморфизмов пространства $\mathbb{R} q$, причем

\end{itemize}



$\Psi_{\gamma \alpha, u}=\Psi_{\gamma \beta, u} \circ \Psi_{\beta \alpha, u} \quad$ при $\quad u \in U_{\alpha} \cap U_{\beta} \cap U_{\gamma}$.



Два хефлигеровских атласа определяют одну и ту же $\mathrm{X}$. с., если они являются частями нек-рого большего хефлигеровского атласа. (Таким образом, Х. с. можно определить как максимальный хефлигеровский атлас.) Если на $X$ задана X. с. $\mathscr{H}$ посредством атласа $\left\{U_{\alpha}\right.$, $\left.\Psi_{\alpha \beta}\right\}$ и $f: Y \longrightarrow X-$ нешрорывное отображение, то атлас $\left\{f^{-1} U_{\alpha}, \Psi_{\alpha \beta}^{\prime}\right\}$, где $\Psi_{\alpha \beta, u}^{\prime}=\Psi_{\alpha \beta, f(u)}$, определяет и нд у ц и р о в анн у іо Х. с. $f^{*} \mathscr{H}$ (она не зависит от кон-



\begin{center}

\includegraphics[max width=\textwidth]{2023_07_09_a01006407529fd695617g-1}

\end{center}





\end{document}